\documentclass[11pt]{article}
\usepackage{amsmath,amssymb,amsthm,amsfonts,hyperref}
\newtheorem{theorem}{Theorem}[section]
\newtheorem{proposition}[theorem]{Proposition}
\newtheorem{lemma}[theorem]{Lemma}
\newtheorem{definition}[theorem]{Definition}
\newtheorem{remark}[theorem]{Remark}

\title{Six Routes to $G(K3 \times E)$: Distinct Constructions, Distinct Invariants}
\author{Raeez Lorgat}
\date{2026-04-22}

\begin{document}
\maketitle

\begin{abstract}
Six routes to the conjectural quantum vertex chiral group
$G(K3 \times E)$ are six distinct constructions,
not six applications of a single functor $\Phi$ to one Calabi--Yau
object. They arise as the six orbits of the Narain orthogonal
group $O(\Lambda_{\mathrm{Nar}})$,
$\Lambda_{\mathrm{Nar}} = \mathrm{II}_{4,20} \oplus \mathrm{II}_{1,1}(-1)$,
acting on primitive Mukai vectors of fixed norm; equivalently, as
the six cosmic-string types of type IIA on $K3 \times E$. Four
of the six produce the canonical spectrum
$\{\kappa_{\mathrm{cat}}, \kappa_{\mathrm{ch}}^{\mathrm{Heis}},
\kappa_{\mathrm{BKM}}, \kappa_{\mathrm{fiber}}\}
= \{0, 3, 5, 24\}$;
the K3 fibre $\kappa_{\mathrm{cat}}(K3) = 2$ is the fibrewise value
and is not the total-space
$\kappa_{\mathrm{cat}}(K3 \times E) = 0$. No symmetric-monoidal
$1$-functor $\Phi \colon \mathrm{CY\text{-}cat}_3 \to \mathrm{Alg}$
has a scalar shadow reproducing all four $\kappa_\bullet$
simultaneously: the four invariants carry four different
multiplicativity types. At the $(\infty,2)$-categorical level the
six routes assemble into a single coherent $6$-simplex
$\sigma \in N_6(\mathcal G_{\mathrm{mn}})$ in the coherent nerve of
the moonshine groupoid, whose seven vertices and fifteen inter-orbit
edges witness the structure. The 6d holomorphic Chern--Simons
theory on $K3 \times E$ gives each route a boundary interpretation:
six different ways to extract the chiral boundary algebra
depending on which cycles of $K3 \times E$ one traces and which
compactification limit one takes.
\end{abstract}

\tableofcontents

\section{What the question asks, and the shape of the answer}
\label{sec:intro}

The symbol $G(K3 \times E)$ denotes the conjectural $(E_1\text{-chiral})$
quantum vertex chiral group associated with the compact Calabi--Yau
threefold $K3 \times E$. Six constructions in the literature compute
something naturally called $G(K3 \times E)$: the Borcherds--Kac--Moody
superalgebra $\mathfrak g_{\Delta_5}$ of Gritsenko--Nikulin
$1995/1997$; the cohomological Hall algebra $\mathrm{CoHA}(K3 \times E)$
of Schiffmann--Vasserot / Kontsevich--Soibelman; the
Maulik--Okounkov stable-envelope quantum group; the
Braverman--Finkelberg--Nakajima Coulomb branch of the
$3$d $\mathcal N = 4$ theory $T[K3 \times E]$; the Kummer-orbifold
free-field vertex algebra; the
Beem--Lemos--Liendo--Peelaers--Rastelli $\Omega$-independent Schur
sector of the $4$d $\mathcal N = 2$ SCFT obtained from $K3 \times E$.
These six constructions do not arise as six applications of one
functor to one object. Each carries a distinct structure; each
computes a distinct invariant; each lands on a distinct
Narain orthogonal orbit of primitive Mukai vectors.

The content of the present note is: (i) a clean identification of
the six constructions, including the two whose status as separate
routes is adversarial (Mukai-lattice vs K3-fibre-rank); (ii) the
Narain orbit stratification that makes the six-fold structure
canonical rather than notational; (iii) the precise statement that
no single symmetric-monoidal $1$-functor reconciles the four
multiplicativity types, and a sharpening that records where a
unifying structure exists ($(\infty,2)$-stratified, not
$1$-categorical); (iv) the 6d holomorphic Chern--Simons
interpretation, in which the six routes are six boundary-reduction
procedures corresponding to six cycle choices; (v) five
attack-heal cycles that test the identifications.

\section{The four invariants}
\label{sec:four-invariants}

Before enumerating constructions, four scalar invariants are
attached to $K3 \times E$, each a different chiral shadow of a
different algebraic structure.

\begin{definition}[Categorical Euler characteristic]
\label{def:kappa-cat}
For a compact Calabi--Yau $d$-fold $X$ with structure sheaf $\mathcal O_X$,
$\kappa_{\mathrm{cat}}(X) := \chi(\mathcal O_X)
= \sum_q (-1)^q h^{0,q}(X)$.
On a product $\kappa_{\mathrm{cat}}$ is K\"unneth-multiplicative:
$\kappa_{\mathrm{cat}}(X \times Y)
= \kappa_{\mathrm{cat}}(X) \cdot \kappa_{\mathrm{cat}}(Y)$.
On $K3$: $\chi(\mathcal O_{K3}) = 1 - 0 + 1 = 2$. On $E$:
$\chi(\mathcal O_E) = 1 - 1 = 0$. On $K3 \times E$:
$\kappa_{\mathrm{cat}}(K3 \times E) = 2 \cdot 0 = 0$.
\end{definition}

\begin{definition}[Heisenberg chiral characteristic]
\label{def:kappa-ch-heis}
$\kappa_{\mathrm{ch}}^{\mathrm{Heis}}(K3 \times E) := 3$
is the degree-$2$ chiral shadow of the $E_1$-chiral algebra
output of $\Phi$ applied to $\mathrm{Perf}(K3 \times E)$; it
equals the Gritsenko weight of the $\mathfrak g_{\Delta_5}$
Cartan plus the coupling variable, signature $(2,1)$, reading
$\mathrm{rk}(\Lambda^{2,1}_{II}) = 3$. Chain-level: five
independent computations agree (chiral de Rham, boundary
chiral algebra of 6d hCS on $K3 \times E$, relative DT
fibre-sewing, BCOV genus-$1$ with corrected
Euler input $\chi^{\mathrm{CY}} = 5$, small $\mathcal N = 4$
$c = 6k$ with $k = 3$).
\end{definition}

\begin{definition}[Borcherds--Kac--Moody weight]
\label{def:kappa-BKM}
For a paramodular Siegel cusp form $\Phi_N$ with Fourier
expansion $\Phi_N(Z) = \sum_{T} c_N(T) e^{2\pi i \operatorname{tr}(TZ)}$,
$\kappa_{\mathrm{BKM}}(\Phi_N) := c_N(0)/2$.
For the Gritsenko--Nikulin $\Delta_5$, $c_1(0) = 10$, so
$\kappa_{\mathrm{BKM}}(\Delta_5) = 5$. The universal formula
$\kappa_{\mathrm{BKM}}(\Phi_N) = c_N(0)/2$ holds unconditionally
for $N \in \{1, 2, 3, 4, 6\}$ (Borcherds $1995$; Gritsenko
$1999$); it is the weight of the Borcherds automorphic product
attached to the lattice $\mathrm{II}_{2,2}(N)$.
\end{definition}

\begin{definition}[Fibre-rank characteristic]
\label{def:kappa-fiber}
$\kappa_{\mathrm{fiber}}(K3 \times E) :=
\chi_{\mathrm{top}}(K3) = 24$. This is the rank of the Mukai
lattice $\Lambda_{\mathrm{Muk}}(K3) = H^\ast(K3, \mathbb Z) \cong
\mathrm{II}_{4,20}$, equivalently the second Chern number
$c_2(K3)$ integrated against the fundamental class, equivalently
$24$ times the topological Euler characteristic of a single
$K3$ fibre. In the relative construction the K3 fibre supplies
$24$ free chiral bosons (lattice VOA $V_{\Lambda_{\mathrm{Muk}}}$).
\end{definition}

\begin{remark}[The four multiplicativity types]
\label{rem:four-mult}
$\kappa_{\mathrm{cat}}$ is K\"unneth-multiplicative;
$\kappa_{\mathrm{fiber}}$ is restriction to a fibre
(neither additive nor multiplicative under products, since
$\kappa_{\mathrm{fiber}}$ depends on the fibration structure,
not only on the total space);
$\kappa_{\mathrm{BKM}}$ is additive under Borcherds tensor product
of the underlying automorphic inputs;
$\kappa_{\mathrm{ch}}^{\mathrm{Heis}}$ is invariant under the
Borcherds tensor product on the chiral side (level is fixed
by the small-$\mathcal N=4$ central charge constraint
$c = 6k$ at $k = 3$).
Four multiplicativity types, four invariants: no single binary
operation on scalars reproduces all four consistently
(Proposition~\ref{prop:no-single-functor}).
\end{remark}

\section{The six routes: construction-by-construction}
\label{sec:six-routes}

Each subsection names (a) the construction; (b) the chiral-side
avatar as an algebraic object $\mathbf A^{(i)}$; (c) the
Narain orbit in which its primitive Mukai vectors live; (d) the
$\kappa_\bullet$ it computes; (e) the status
(proved / conjectural).

\subsection{Route R$1$: Mukai-lattice real-root orbit}
\label{subsec:R1-mukai}

\emph{Construction.} The Mukai lattice $\Lambda_{\mathrm{Muk}}
= H^\ast(K3, \mathbb Z)
\cong 2U \oplus 2E_8(-1) \oplus U$, equivalently the even unimodular
lattice $\mathrm{II}_{4,20}$ of signature $(4, 20)$, rank $24$.
The associated chiral algebra is the Mukai-lattice vertex
algebra $V_{\Lambda_{\mathrm{Muk}}}$ in the signature-enhanced
Heisenberg package; on the K3 fibre this reads the $24$ free
chiral bosons organised by $\Lambda_{\mathrm{Muk}}$.
\emph{Scalar shadow.} At the fibre level,
$\kappa_{\mathrm{cat}}(K3) = \chi(\mathcal O_{K3}) = 2$, the
Hodge supertrace $(1 + 1)$ on $K3$. The $2$ is the K3 fibre
$\kappa_{\mathrm{cat}}$, not the total-space $\kappa_{\mathrm{cat}}$.
\emph{Orbit.} Real-root orbit, discriminant $D < 0$, gcd $=1$:
primitive vectors $v$ with $\langle v, v \rangle = -2$ in the
$\mathrm{II}_{3,19}$ sublattice.
\emph{Chiral-side avatar.} The $K3$-vertex-algebra truncation
$V_{\mathrm{II}_{3,19}}$ with $\kappa_{\mathrm{ch}}^{\mathrm{cat}}
= 2$.

\subsection{Route R$2$: Heisenberg $E_1$-chiral orbit}
\label{subsec:R2-heis}

\emph{Construction.} The $E_1$-chiral algebra
$\mathbf H_{\Delta_5}$ produced by the CY-to-chiral functor $\Phi$
on $\mathrm{Perf}(K3 \times E)$ at $d = 3$. This is the unique
route involving $\Phi$; the remaining five produce their outputs
by independent means. The construction passes through
factorisation homology
$\int_{K3 \times E} (-)$ of a Calabi--Yau factorisation algebra,
with dimension $d = 3$ forcing the output to carry $E_1$-chiral
rather than $E_2$-chiral structure.
\emph{Scalar shadow.}
$\kappa_{\mathrm{ch}}^{\mathrm{Heis}}(K3 \times E) = 3$
(Definition~\ref{def:kappa-ch-heis}).
\emph{Orbit.} Heisenberg orbit, $D < 0$, gcd $>1$: imprimitive
Mukai vectors along the $\mathrm{II}_{2,2}$ sublattice generated
by the Cartan of $\mathfrak g_{\Delta_5}$ plus the elliptic
coupling $z$.
\emph{Chiral-side avatar.}
$\mathbf H_{\Delta_5}$ as Kontsevich--Soibelman
$\mathrm{CoHA}(K3 \times E)$-module, character identity
$\chi(\mathbf H_{\Delta_5}) = Z_{K3 \times E}^{\mathrm{DT}}
/ (\text{volume factor})
= 1 / \Delta_5(Z)^2
= 1 / \Phi_{10}(Z)$.

\subsection{Route R$3$: Light-like Igusa squaring orbit ($\Phi_{10} = \Delta_5^2$)}
\label{subsec:R3-igusa}

\emph{Construction.} The Igusa cusp form
$\Phi_{10} \in S_{10}(\mathrm{Sp}_4(\mathbb Z))$ is the unique
weight-$10$ Siegel cusp form with trivial character; it factors
as $\Phi_{10} = \Delta_5^2$ with $\Delta_5 \in
S_5(\mathrm{Sp}_4(\mathbb Z), \nu_{\Delta_5})$ the
Gritsenko--Nikulin paramodular form of weight $5$ with
nebentypus. The DT partition function on $K3 \times E$ reads
$Z_{K3 \times E}^{\mathrm{DT}} = C / \Phi_{10} = C / \Delta_5^2$
(Oberdieck $2018$). The construction of $\mathfrak g_{\Delta_5}$
by Borcherds lift of the weak Jacobi form $\phi_{0,1}$
(the K3 elliptic genus / $2$) is the additive-lift side;
the multiplicative-lift side produces $\Delta_5$ itself
as the denominator identity.
\emph{Scalar shadow.}
$\kappa_{\mathrm{BKM}}(\Phi_{10}) = c_{10}(0)/2 = 10$; not in
the canonical K3 $\times$ E tetrad but present as the
duality-orbit squaring of R$4$. On the light-like ray, the
shadow is $\eta^9$-multiplicity, $\tau(a) = 9$.
\emph{Orbit.} Light-like orbit, $D = 0$: primitive null vectors
in $\mathrm{II}_{4,20} \oplus \mathrm{II}_{1,1}(-1)$.
\emph{Chiral-side avatar.} Borcherds tensor square
$V_{\Phi_{10}} = V_{\Delta_5} \otimes V_{\Delta_5}$, weight $10$.

\subsection{Route R$4$: BKM paramodular orbit ($\Phi_N$, $N \in \{1,2,3,4,6\}$)}
\label{subsec:R4-BKM}

\emph{Construction.} The Borcherds--Kac--Moody superalgebra
$\mathfrak g_{\Delta_5}$ is the $N = 1$ (untwisted) member of
the paramodular family $\{\Phi_N\}_{N \in \{1, 2, 3, 4, 6\}}$
of Gritsenko--Nikulin. Each $\Phi_N$ is a Borcherds
automorphic product on the paramodular group
$K(N) = \mathrm{II}_{2,2}(N) \cdot \mathrm{Sp}_4$; its Fourier
coefficients compute root multiplicities
$\mathrm{mult}(\alpha) = c_N(4nm - \ell^2)$.
\emph{Scalar shadow.}
$\kappa_{\mathrm{BKM}}(\Phi_N) = c_N(0)/2$ for $N \in \{1, 2, 3, 4, 6\}$
(Borcherds $1995$; Gritsenko $1999$); for $N = 1$,
$c_1(0) = 10$, $\kappa_{\mathrm{BKM}}(\Delta_5) = 5$.
\emph{Orbit.} Imaginary orbit, $D > 0$, gcd $= 1$: positive-norm
primitive vectors.
\emph{Chiral-side avatar.} The denominator-formula chiral algebra
$V_0(\mathfrak g_{\Delta_5})$ reading the Borcherds product
$\Delta_5(Z) = q_1^{1/2} q_2^{-1/2} q_3^{1/2}
\prod_{(m,n,\ell) > 0} (1 - q_1^m q_2^\ell q_3^n)^{c_1(4mn - \ell^2)}$.

\subsection{Route R$5$: Monster / Niemeier / $N$-twisted orbit ($\gcd = N$)}
\label{subsec:R5-Monster}

\emph{Construction.} The Borcherds Monster Lie algebra
$\mathfrak m \subset V^\natural$ and its Niemeier siblings,
with root lattice $\mathrm{II}_{25,1}$, realising the
Conway--Norton moonshine module. On $K3 \times E$ the
Monster-orbit route is the Borcherds $1992$ embedding
$\mathrm{II}_{3,19} \hookrightarrow \mathrm{II}_{25,1}$
composed with the Fake-Monster/Monster Lie superalgebra
production rule; it is the route through which the Monster
twinings act on the K3 elliptic genus (EOT moonshine,
Mathieu $M_{24}$ twining classes).
\emph{Scalar shadow.}
$\kappa_{\mathrm{ch}}(V^\natural) = 12$, distinct from the
rank $24$ of the Leech-lattice VOA $V_\Lambda$; the
half-shadow $12 = 24/2$ encodes the order-$2$ Golay-code
twist. For $N$-twisted cases, $\kappa_\bullet$ reads off
the $N$-th Niemeier orbit discriminant.
\emph{Orbit.} $N$-twisted orbit, gcd $= N$: imprimitive Mukai
vectors of fixed denominator.
\emph{Chiral-side avatar.} $V^\natural$ at $N = 1$, Niemeier
lattice VOAs $V_{N_i}$ at higher $N$.

\subsection{Route R$6$: K3 fibre-rank orbit}
\label{subsec:R6-fiber}

\emph{Construction.} The relative chiral algebra construction
$\pi \colon K3 \times E \to E$ produces the fibration
$K3$-fibre-by-$K3$-fibre, and sews the fibres along $E$ via the
DMVV / elliptic-genus formula. Each $K3$ fibre contributes a
boundary algebra with $\chi_{\mathrm{top}}(K3) = 24$ zero modes;
the $24$ reads the Mukai lattice rank and also the anomaly
coefficient of the $K3$ sigma model.
\emph{Scalar shadow.}
$\kappa_{\mathrm{fiber}}(K3 \times E) = 24$
(Definition~\ref{def:kappa-fiber}).
\emph{Orbit.} Fibre-rank orbit: the anti-diagonal embedding
$\mathrm{II}_{3,19} \hookrightarrow \mathrm{II}_{4,20}
\oplus \mathrm{II}_{1,1}(-1)$ selecting the K3-fibre summand.
\emph{Chiral-side avatar.} The relative vertex algebra
$\pi_\ast V_{\Lambda_{\mathrm{Muk}}}$ on $E$, equivalently the
$24$-boson lattice VOA $V_{\Lambda_{\mathrm{Muk}}}$ evaluated on
the $K3$-fibre, with Euler contribution $24$ sewn along the
elliptic fibre.

\subsection{BLLPR as a sixth Route (Schur-sector reading of R$2$)}
\label{subsec:BLLPR-sixth}

A sixth construction of $G(K3 \times E)$ available in the
literature is the Beem--Lemos--Liendo--Peelaers--Rastelli
Schur sector: the $4$d $\mathcal N = 2$ SCFT $T[K3 \times E]$
obtained by compactifying the $6$d $(2,0)$ theory on the
Calabi--Yau threefold; its Schur sector carries a VOA
structure by the BLLPR theorem; for $K3 \times E$ this VOA is
the small $\mathcal N = 4$ SCA at central charge $c = 6k$.
\emph{Feature.} The BLLPR construction is $\Omega$-independent:
it uses no $\Omega$-background localisation.
\emph{Reading.} At the chiral-side scalar shadow,
the BLLPR VOA lands on $k = \kappa_{\mathrm{ch}}^{\mathrm{Heis}} = 3$
(sigma-model reading), equivalently $k = 2$ (abelian-free-field
reading). This is R$2$'s chiral-side object reproduced by a
route through holography of the $4$d $\mathcal N = 2$ index, not
an independent orbit. In the adversarial sense
(\S\ref{subsec:attack-1}), BLLPR is a sharpening of R$2$, not
a seventh route, but the sharpening is nontrivial: the
Schur-sector description makes the $\mathcal N = 4$ SCA
structure manifest in a way that the factorisation-homology
description of $\Phi$ does not.

\section{The Narain orthogonal orbit structure}
\label{sec:orbits}

\begin{definition}[Narain orbit decomposition]
\label{def:narain-orbits}
Let $\Lambda_{\mathrm{Nar}} = \mathrm{II}_{4,20} \oplus
\mathrm{II}_{1,1}(-1)$ of signature $(5, 21)$, and let
$G_{\mathrm{Nar}} = O(\Lambda_{\mathrm{Nar}})$ act on
primitive lattice vectors of fixed Mukai norm
$\mathrm{Muk}(v) := \langle v, v \rangle_{\Lambda_{\mathrm{Nar}}}
\in \mathbb Z$.
For each orbit type $i \in \{\mathrm I, \mathrm{II}, \mathrm{III},
\mathrm{IV}, \mathrm V, \mathrm{VI}\}$ the characteristic data
$(\text{discriminant } D, \gcd, \text{sublattice})$ are:
\[
\begin{array}{lcccl}
\text{Orbit} & D & \gcd & \text{Sublattice} & \kappa_\bullet \\
\hline
\mathrm I \ (\text{Mukai real-root}) & <0 & 1 & \mathrm{II}_{3,19} & \kappa_{\mathrm{cat}}(K3) = 2 \\
\mathrm{II} \ (\text{Heisenberg}) & <0 & >1 & \mathrm{II}_{2,2} & \kappa_{\mathrm{ch}}^{\mathrm{Heis}} = 3 \\
\mathrm{III} \ (\text{light-like}) & 0 & 1 & \mathbb Z \cdot \delta & \eta^9, \tau(a) = 9 \\
\mathrm{IV} \ (\text{imaginary/BKM}) & >0 & 1 & \mathrm{II}_{2,3} & \kappa_{\mathrm{BKM}} = 5 \\
\mathrm V \ ($N$-twisted) & \text{mixed} & N & \mathrm{II}_{25,1}(N) & \Phi_N\text{-family} \\
\mathrm{VI} \ (\text{fibre-rank}) & \text{mixed} & 1 & \mathrm{II}_{4,20} & \kappa_{\mathrm{fiber}} = 24.
\end{array}
\]
\end{definition}

\begin{theorem}[Orbit--Route correspondence]
\label{thm:orbit-route}
The six routes R$1, \ldots,$ R$6$ of \S\ref{sec:six-routes} are in
canonical bijection with the six Narain orbits
$(\mathrm I), \ldots, (\mathrm{VI})$ of
Definition~\ref{def:narain-orbits}:
R$i$ is the chiral-side avatar attached to the sublattice of orbit
$(i)$ via the corresponding Borcherds/Gritsenko/Schiffmann--Vasserot/\-
Maulik--Okounkov/Niemeier/Mukai construction.
\end{theorem}

\begin{proof}[Proof sketch]
Pass through the Narain moduli stratification of
type-IIA $K3 \times E$: primitive BPS states are Mukai vectors
in $\Lambda_{\mathrm{Nar}}$, and the BPS-spectrum generating function
$Z_{\mathrm{BPS}}(K3 \times E) = 1/\Phi_{10}$ (Oberdieck $2018$) has
Fourier coefficients stratified by orbit type. Each stratum $i$
has its own automorphic-lift construction producing a distinct
chiral algebra $\mathbf A^{(i)}$: $(\mathrm I)$ via Mukai real-root
VOA, $(\mathrm{II})$ via Heisenberg averaging, $(\mathrm{III})$ via
Borcherds tensor squaring, $(\mathrm{IV})$ via Gritsenko paramodular
lift, $(\mathrm V)$ via Niemeier lattice VOA + Conway twinings,
$(\mathrm{VI})$ via relative fibre DMVV sewing. The bijection is
forced by the orbit-stabiliser decomposition of
$O(\Lambda_{\mathrm{Nar}}) \curvearrowright
\mathrm{Prim}_{(2,26)}$; see
\S\ref{sec:six-simplex} for the homotopy-coherence structure.
\end{proof}

\section{The adversarial question: are six routes really six?}
\label{sec:adversarial}

\subsection{Attack 1: Mukai vs fibre-rank (R$1$ vs R$6$ collapse?)}
\label{subsec:attack-1}

\emph{Attack.} R$1$ is the Mukai lattice $\Lambda_{\mathrm{Muk}}
\cong \mathrm{II}_{4,20}$ of rank $24$ signature $(4,20)$; R$6$ is
the fibre-rank $\chi_{\mathrm{top}}(K3) = 24$, with K3 Mukai lattice
rank $24$. Are R$1$ and R$6$ the same construction under different
names?

\emph{Heal.} No: the sublattice in which primitive vectors live
differs. R$1$ uses the $\mathrm{II}_{3,19}$ sublattice
(Picard real-root real cone on primitive Picard elements,
$D < 0$, gcd $= 1$), producing the K3-fibrewise
$\kappa_{\mathrm{cat}}(K3) = 2$ via the Hodge supertrace
$1 - 0 + 1$ on the Dolbeault complex of the fibre. R$6$ uses the
full Mukai lattice $\mathrm{II}_{4,20}$ together with the
coupling cycle $\mathrm{II}_{1,1}(-1)$, producing the
fibrewise $\chi_{\mathrm{top}}(K3) = 24$ via the topological
Euler characteristic of the fibre. The two sublattices are
different: $\mathrm{II}_{3,19} \subset \mathrm{II}_{4,20}$ with
quotient $U$ (the hyperbolic summand corresponding to
momentum/winding on a null direction). At the chiral-side
avatar level, R$1$ is the truncated-Mukai $V_{\mathrm{II}_{3,19}}$
and R$6$ is the full $V_{\Lambda_{\mathrm{Muk}}}$. The difference
is one hyperbolic summand, and the corresponding chiral-side
correction is the passage from the $22$-generator Heisenberg
presentation of $Y^{\mathrm{Heis}}_\hbar(\Lambda_{K3})$ to the
$24$-generator extension including the null pair. Two distinct
constructions, producing two distinct numerical invariants:
$\kappa_{\mathrm{cat}}(K3) = 2$ vs $\kappa_{\mathrm{fiber}} = 24$.

\subsection{Attack 2: Igusa squaring vs BKM (R$3$ vs R$4$ collapse?)}
\label{subsec:attack-2}

\emph{Attack.} $\Phi_{10} = \Delta_5^2$ by the Gritsenko--Nikulin
squaring; so R$3$ (Igusa $\Phi_{10}$) is R$4$ (BKM $\Delta_5$)
raised to the second power. Are they the same construction?

\emph{Heal.} No: the scalar shadow differs,
$\kappa_{\mathrm{BKM}}(\Phi_{10}) = c_{10}(0)/2 = 10$ while
$\kappa_{\mathrm{BKM}}(\Delta_5) = c_1(0)/2 = 5$; and the
orbit differs, R$3$ living on the light-like ray $D = 0$
with $\eta^9$-multiplicity datum $\tau(a) = 9$, R$4$ on the
imaginary orbit $D > 0$ with BKM root multiplicity
$c_1(4mn - \ell^2)$. The light-like $\eta^9$ datum is
$\tau(a) = 9 = c_1(0) - 1 = 2\kappa_{\mathrm{BKM}} - 1$, a
shifted-shadow of R$4$ by one unit, consistent with the
Atkin--Lehner squaring map on automorphic inputs but not
reducing the two routes to one.

\subsection{Attack 3: Monster vs BKM (R$5$ vs R$4$ collapse?)}
\label{subsec:attack-3}

\emph{Attack.} The Monster Lie algebra $\mathfrak m$ is a BKM
superalgebra (Borcherds $1992$); so is $\mathfrak g_{\Delta_5}$.
Both embed into a common Borcherds framework. Why are they
different routes?

\emph{Heal.} Different root lattices: $\mathfrak m$ has
$\mathrm{II}_{25,1}$, $\mathfrak g_{\Delta_5}$ has
$\Lambda^{2,1}_{II}$. Different automorphic inputs: $\mathfrak m$
uses the Borcherds $1992$ denominator formula
$j(\tau) - j(\sigma) = p^{-1} \prod_{m, n > 0}
(1 - p^m q^n)^{c(mn)}$ with $c(n)$ the $j$-function
coefficients; $\mathfrak g_{\Delta_5}$ uses the paramodular
Borcherds product for $\Delta_5$ with $\phi_{0,1}$ input. The
embedding $\mathrm{II}_{3,19} \hookrightarrow \mathrm{II}_{25,1}$
is one of the $24$ Niemeier extensions and links the two routes
at the $1$-morphism level (edge $e_{45}$ in
\S\ref{sec:six-simplex}), but does not identify them.

\subsection{Attack 4: CoHA vs MO vs BFN (are R$2$'s equivalents
multiple?)}
\label{subsec:attack-4}

\emph{Attack.} The Schiffmann--Vasserot CoHA, the Maulik--Okounkov
stable-envelope construction, the Braverman--Finkelberg--Nakajima
Coulomb branch, and the Kummer-orbifold free-field VOA all produce
algebras of relevance to $G(K3 \times E)$ and all are labelled
``R$2$-like''. Are they $4$ separate constructions, making the
true route count $9$?

\emph{Heal.} They are $1$-morphisms between R$2$ and its avatars
within the Heisenberg orbit $(\mathrm{II})$, not independent
routes. Concretely: Schiffmann--Vasserot CoHA is the positive-half
functor $\mathrm{CoHA}(K3) = Y^+_\hbar(\Lambda_{K3})$ (Theorem,
Schiffmann--Vasserot $2017$). Maulik--Okounkov stable envelopes
produce the associated $R$-matrix via stable bases on
$\mathrm{Hilb}^n(K3)$. BFN Coulomb branches produce
$Y_\hbar(\Lambda_{K3})$ as a deformation quantisation of
$\mathcal M_C(T[K3])$. Kummer orbifold gives a free-field
model at the orbifold point $T^4/\mathbb Z_2$. All four land in
the same Heisenberg orbit $(\mathrm{II})$ and produce the same
$\kappa_{\mathrm{ch}}^{\mathrm{Heis}} = 3$; they are related by
equivalences of stable $\infty$-categories (proved where the
chain-level computation aligns; conjectural for the
Coulomb-branch-to-CoHA comparison at full K3). The six-route
structure counts $O(\Lambda_{\mathrm{Nar}})$-orbits, not
constructions-up-to-equivalence; each orbit admits multiple
equivalent constructions without multiplying the route count.

\subsection{Attack 5: can a single functor unify the four invariants?}
\label{subsec:attack-5}

\emph{Attack.} Six constructions compute four invariants;
the surplus suggests that a unifying functor exists that produces
a single algebraic output carrying all four scalar shadows.

\emph{Heal.} No, at the $1$-categorical level.

\begin{proposition}[No single-target $1$-functor unification]
\label{prop:no-single-functor}
No symmetric-monoidal functor
$\Phi^{\mathrm{gen}} \colon \mathrm{CY\text{-}cat}_3 \to
\mathrm{Alg}$ into a single $1$-category of algebras has a scalar
shadow simultaneously computing all four of
$\kappa_{\mathrm{cat}}$, $\kappa_{\mathrm{ch}}^{\mathrm{Heis}}$,
$\kappa_{\mathrm{BKM}}$, $\kappa_{\mathrm{fiber}}$ on $K3 \times E$.
\end{proposition}

\begin{proof}
Suppose $\Phi^{\mathrm{gen}}$ symmetric-monoidal with scalar shadow
$f$ satisfying $f(A \boxtimes A') = f(A) \ast f(A')$ for a fixed
binary operation $\ast$. Apply to $A = \mathcal A_{K3}$, $A' =
\mathcal A_E$. If $f = \kappa_{\mathrm{cat}}$ then $\ast = \cdot$
and $\kappa_{\mathrm{cat}}(K3) \cdot \kappa_{\mathrm{cat}}(E)
= 2 \cdot 0 = 0 = \kappa_{\mathrm{cat}}(K3 \times E)$: consistent.
If $f = \kappa_{\mathrm{fiber}}$ then $24 = 24 \ast
\kappa_{\mathrm{fiber}}(E)$ must hold, but $\kappa_{\mathrm{fiber}}$
is not defined on a generic elliptic factor (no fibration
structure), and no associative $\ast$ with neutral element on
$\mathbb Z$ returns $24$ as the product of $24$ and an undefined
input; inconsistent. If $f = \kappa_{\mathrm{BKM}}$ then $5
= \kappa_{\mathrm{BKM}}(K3) \ast \kappa_{\mathrm{BKM}}(E)$, but
$\kappa_{\mathrm{BKM}}$ is only defined on outputs of Borcherds
lifts and not on individual CY factors. No single binary
operation reconciles all four; contradiction.
\end{proof}

\begin{theorem}[Stratified unifying functor at $(\infty, 2)$]
\label{thm:stratified-unification}
There exists a unifying functor
\[
\Phi^{\mathrm{gen}} \colon \mathrm{CY\text{-}cat}_3
\longrightarrow
\mathrm{Strat}_{(\infty, 2)}\bigl(\mathrm{UDual}\bigr)
\]
into the $(\infty, 2)$-category of U-duality-stratified outputs,
defined by
\[
\Phi^{\mathrm{gen}}(\mathcal A_X) := \bigl\{
\Phi^{(\mathrm I)}(\mathcal A_X), \ldots,
\Phi^{(\mathrm{VI})}(\mathcal A_X) \bigr\},
\]
with $\Phi^{(i)}$ the restriction of the BPS root-space sheaf to the
$i$-th orbit. For $X = K3 \times E$ the output recovers the full
spectrum $\{0, 3, 5, 24\}$ via the orbit-collision map
$\mathsf{coll} \colon \{\mathrm I, \ldots, \mathrm{VI}\}
\twoheadrightarrow \{\kappa_{\mathrm{cat}},
\kappa_{\mathrm{ch}}^{\mathrm{Heis}},
\kappa_{\mathrm{BKM}}, \kappa_{\mathrm{fiber}}\}$,
with orbits $(\mathrm I)$ and $(\mathrm{VI})$ distinguished by
K3-fibre vs fibre-rank, orbits $(\mathrm{III})$ and $(\mathrm V)$
organised by the $\Phi_N$-paramodular family, and orbits
$(\mathrm{II})$ and $(\mathrm{IV})$ the
$\kappa_{\mathrm{ch}}^{\mathrm{Heis}}$ and
$\kappa_{\mathrm{BKM}}$ poles.
\end{theorem}

\begin{proof}[Sketch]
$\Lambda_{\mathrm{Nar}}$ carries the $O(\Lambda_{\mathrm{Nar}})$
U-duality action; primitive vectors of fixed Mukai norm fall into
six orbit types by Eichler's classification. The BPS root-space
sheaf $\mathcal R_{\mathrm{BPS}} \to
\mathrm{Narain}/O(\Lambda_{\mathrm{Nar}})$ restricts to each orbit,
yielding a distinct algebraic invariant; the $(\infty, 2)$-groupoid
structure records Kontsevich--Soibelman BPS wall-crossing
$1$-morphisms and U-duality isotopies as $2$-morphisms.
\end{proof}

\section{The $6$-simplex in the moonshine nerve}
\label{sec:six-simplex}

The six routes are not a list; they are the edges out of a
basepoint in a $6$-simplex whose homotopy coherence records the
fifteen inter-route comparisons and the $35$ triangle
factorisations.

\begin{definition}[The seven vertices]
\label{def:seven-vertices}
Let $\mathcal G_{\mathrm{mn}}$ be the $(\infty, 1)$-groupoid of
moonshine automorphic data, with objects
$(L, F, \Psi, \iota)$ comprising an even lattice $L$, a Siegel
automorphic form $F$ on $L$, an automorphic data
$\Psi$, and a Fourier expansion $\iota$. Define seven vertices
$(v_0, v_1, \ldots, v_6) \in \mathrm{Ob}(\mathcal G_{\mathrm{mn}})$:
\begin{align*}
v_0 &= (\mathrm{II}_{3,19}, \Delta_5, \Psi_{\mathrm{Grit}},
\iota_{\Delta_5}) \quad \text{(basepoint)}, \\
v_1 &= (\mathrm{II}_{3,19}, \Theta_{K3}, \Psi^{\mathrm{Muk}},
\iota^{(\mathrm I)})  \quad \text{(R$1$, Mukai real-root)}, \\
v_2 &= (\mathrm{II}_{2,2}, \eta^{24}, \Psi^{\mathrm{Heis}},
\iota^{(\mathrm{II})})  \quad \text{(R$2$, Heisenberg)}, \\
v_3 &= (2U \oplus 2E_8(-1), \Delta_{10}, \Psi^{(2)}_{\mathrm{Grit}},
\iota^{(\mathrm{III})})  \quad \text{(R$3$, light-like)}, \\
v_4 &= (\mathrm{II}_{2,3}, \Phi_N, \Psi^{\mathrm{para}},
\iota^{(\mathrm{IV})})  \quad \text{(R$4$, imaginary/BKM)}, \\
v_5 &= (\mathrm{II}_{25,1}, \Delta_{\mathrm{Mon}}, \Psi^{\mathrm{Bor}92},
\iota^{(\mathrm V)})  \quad \text{(R$5$, $N$-twisted Monster)}, \\
v_6 &= (\mathrm{II}_{4,20} \oplus \mathrm{II}_{1,1}(-1),
Z_{K3}/\eta^{24}, \Psi^{\mathrm{Mat}},
\iota^{(\mathrm{VI})})  \quad \text{(R$6$, fibre-rank)}.
\end{align*}
\end{definition}

\begin{theorem}[The $6$-simplex $\sigma \in N_6(\mathcal G_{\mathrm{mn}})$]
\label{thm:six-simplex}
There is a coherent $6$-simplex
$\sigma \colon \Delta^6 \to N(\mathcal G_{\mathrm{mn}})$ with
vertices of Definition~\ref{def:seven-vertices} and edges
$e_{ij} = \sigma|_{[i, j]}$, $0 \leq i < j \leq 6$:
\begin{itemize}
\item $e_{0i} = r_i$, the $i$-th route $1$-morphism,
$i = 1, \ldots, 6$;
\item $e_{ij}$ for $1 \leq i < j \leq 6$ is the inter-orbit
comparison $c_{ij} \colon v_i \to v_j$ given by the U-duality
lattice embedding $L_i \hookrightarrow L_j$ composed with
Borcherds-multiplicative structure.
\end{itemize}
For each triangle $[i, j, k]$, $0 \leq i < j < k \leq 6$, the
$2$-simplex $\sigma|_{[i,j,k]}$ is the explicit homotopy
$e_{jk} \circ e_{ij} \simeq e_{ik}$. All $\binom{7}{3} = 35$
triangles commute up to homotopy; all $\binom{7}{4} = 35$
tetrahedra, $\binom{7}{5} = 21$ $4$-simplices,
$\binom{7}{6} = 7$ $5$-simplices cohere; the single $6$-simplex
$\sigma$ is the simultaneous witness. The seven codim-$1$ faces
$d_i \sigma$, $0 \leq i \leq 6$, each delete one vertex and
record the corresponding route-deletion content.
\end{theorem}

\begin{remark}[The inter-orbit comparisons]
\label{rem:intercomparisons}
The fifteen inter-orbit edges $c_{ij}$ are concrete:
$c_{13}$ is the Borcherds tensor square
$\Delta_5 \mapsto \Delta_{10}$; $c_{12}$ is the Mukai-to-Heisenberg
projection $\mathrm{II}_{3,19} \to \mathrm{II}_{2,2}$ along the
hyperbolic summand; $c_{14}$ is the Gritsenko additive lift
$\mathrm{II}_{3,19} \hookrightarrow \mathrm{II}_{2,3}$;
$c_{15}$ is the Leech-no-roots/K3-no-$(-2)$ embedding
$\mathrm{II}_{3,19} \hookrightarrow \mathrm{II}_{25,1}$; $c_{16}$
is the fibration-period map
$\mathrm{II}_{3,19} \to \mathrm{II}_{4,20} \oplus
\mathrm{II}_{1,1}(-1)$; remaining $c_{ij}$ with $2 \leq i < j$
are composites through $v_0$.
\end{remark}

\section{6d holomorphic Chern--Simons on $K3 \times E$: six boundary
reductions}
\label{sec:6dhCS}

\begin{definition}[6d hCS on a CY$_3$]
\label{def:6dhCS}
For a Calabi--Yau $3$-fold $X$ with holomorphic volume form $\Omega$,
the Costello $6$d holomorphic Chern--Simons theory at gauge algebra
$\mathfrak g$ has action
\[
S_{\mathrm{hCS}}[A]
= \int_{X} \Omega \wedge \bigl(\tfrac{1}{2} A \,\bar\partial A
+ \tfrac{1}{3} A \, [A, A]\bigr),
\]
where $A \in \Omega^{0, \ast}(X, \mathfrak g)$ is a $\mathfrak g$-valued
Dolbeault form. The theory is topological in the antiholomorphic
direction and holomorphic in the holomorphic direction;
factorisation algebras of observables on $X$ are $E_3$-algebras in
the antiholomorphic directions but carry additional
$\bar\partial$-structure in the holomorphic directions. The boundary
algebra is extracted by specialising to a codimension-$2$ cycle
$\Sigma \subset X$.
\end{definition}

\begin{theorem}[Six boundary reductions of 6d hCS on $K3 \times E$]
\label{thm:six-reductions}
Let $X = K3 \times E$. Corresponding to the six Narain orbits
$(\mathrm I), \ldots, (\mathrm{VI})$ of
Definition~\ref{def:narain-orbits} there are six boundary-reduction
procedures of the 6d hCS theory on $X$, each producing a distinct
chiral-side avatar:
\begin{enumerate}[label=(B\arabic*)]
\item (B$1$) Reduction to $\mathrm{Pic}(K3)$-cycles: traces primitive
$(-2)$-classes of the K3 fibre; boundary is the Mukai real-root
VOA $V_{\mathrm{II}_{3,19}}$ with $\kappa_{\mathrm{cat}}(K3) = 2$.
\item (B$2$) Reduction along the $E$-cycle: traces the Heisenberg
$\hbar$-deformation along the elliptic base; boundary is the
$E_1$-chiral algebra $\mathbf H_{\Delta_5}$ with
$\kappa_{\mathrm{ch}}^{\mathrm{Heis}} = 3$. This is the boundary
dual to $\Phi$ on $\mathrm{Perf}(K3 \times E)$.
\item (B$3$) Squared doubling: takes two copies of B$2$ and contracts
along the diagonal Borcherds tensor square; boundary is
$V_{\Phi_{10}}$ with Igusa weight $10$.
\item (B$4$) Reduction via paramodular period integration:
integrates $\Omega \wedge A^3$ against the paramodular form $\Phi_N$;
boundary is $V_0(\mathfrak g_{\Delta_5})$ with
$\kappa_{\mathrm{BKM}} = 5$ at $N = 1$.
\item (B$5$) Reduction via Niemeier lattice embedding: traces
$(-2)$-roots of a Niemeier lattice via
$\mathrm{II}_{3,19} \hookrightarrow \mathrm{II}_{25,1}$; boundary is
the $V^\natural$-module at $N = 1$, Niemeier VOAs at higher $N$.
\item (B$6$) Relative reduction over $E$: fibre-by-fibre pushforward
$\pi_\ast$ of 6d hCS on $K3 \times E \to E$; boundary is the
$24$-boson lattice VOA $V_{\Lambda_{\mathrm{Muk}}}$ on $E$ with
$\kappa_{\mathrm{fiber}} = 24$.
\end{enumerate}
The six boundary-reduction procedures are related by the homotopy
coherence of the $6$-simplex $\sigma \in N_6(\mathcal G_{\mathrm{mn}})$
of Theorem~\ref{thm:six-simplex}.
\end{theorem}

\begin{proof}[Proof sketch]
The six procedures are the six boundary conditions consistent with
the $\mathrm{CY}_3$ structure on $X$ and the action of the Narain
U-duality on the $\mathrm{II}_{4,20} \oplus \mathrm{II}_{1,1}(-1)$
charge lattice. For each orbit $i$, the corresponding boundary
cycle is the primitive lattice vector representative modulo
$O(\Lambda_{\mathrm{Nar}})$-stabiliser; the boundary algebra is
extracted by evaluating the 6d hCS partition function along this
cycle and reading the factorisation-algebra restriction to the
boundary. That the six boundary algebras assemble into a coherent
homotopy-commutative diagram is Theorem~\ref{thm:six-simplex}.
\end{proof}

\begin{remark}[$E_n$-structure on the boundary]
\label{rem:En-structure}
At $d = 3$ the 6d hCS output carries $E_1$-chiral structure on the
boundary algebra; the $E_2$-chiral structure lives one level up on
$\mathcal Z(\mathrm{Rep}^{E_1}(\mathbf H_{\Delta_5}))$. The ZTE
failure at $O(\kappa_{\mathrm{ch}}^2)$ (the tetrahedron equation obstruction
at second order) is the structural witness that the
$6$d theory is not the dimensional upgrade of its $3$d/$5$d
siblings; it is genuinely new. Boundary reductions B$1$--B$6$
each read $E_1$ on their respective boundary algebras; the
$E_2$-centre arises only through assembly of the six.
\end{remark}

\section{Cross-consistency with cache facts}
\label{sec:consistency}

\begin{proposition}[Consistency summary]
\label{prop:consistency}
The identification of six routes R$1$--R$6$ of \S\ref{sec:six-routes},
their Narain orbits of \S\ref{sec:orbits}, and their boundary
reductions of \S\ref{sec:6dhCS} satisfy:
\begin{enumerate}
\item $\kappa_{\mathrm{cat}}(K3 \times E) = 0$ on the total space by
K\"unneth multiplicativity ($2 \cdot 0 = 0$), distinct from the
K3 fibre $\kappa_{\mathrm{cat}}(K3) = 2$ delivered by R$1$.
\item The six routes are six different constructions, not six
applications of $\Phi$ to one object; only R$2$ involves $\Phi$
(\S\ref{subsec:R2-heis}).
\item $\kappa_{\mathrm{BKM}}(\Phi_N) = c_N(0)/2$ for
$N \in \{1, 2, 3, 4, 6\}$; the $N = 1$ entry
$\kappa_{\mathrm{BKM}}(\Delta_5) = 5$ at R$4$
reproduces the Gritsenko--Nikulin paramodular Fourier coefficient.
$\mathrm{CoHA}(\mathbb C^3) = Y^+$ (positive half) (not
$\mathcal W_{1 + \infty}$), so the CoHA$\to$R$2$ avatar is the
positive-half of the K3 Yangian, not the full affine $\mathcal W$.
\item At $d \geq 3$, the $A = \mathbf H_{\Delta_5}$ produced by R$2$
is $E_1$; the $E_2$-chiral structure lives on
$\mathcal Z(\mathrm{Rep}(A))$, not on $A$.
Class $\mathbf M$ shadow tower for R$2$ has
$E_3$-bar $= 6^g$ at genus-$g$ cohomology
(for the class-$\mathsf M$ shadow of the K3 sigma model).
\item The false identity
$\kappa_{\mathrm{BKM}} = \kappa_{\mathrm{ch}}^{\mathrm{Heis}}
+ \chi(\mathcal O_{\mathrm{fiber}})$
fails at every $N$ (at $N = 1$: $5 \neq 0 + 0 = 0$).
The correct identity is
$\kappa_{\mathrm{BKM}}(\Phi_N) = c_N(0)/2$, a paramodular
Fourier constant-term identity with no decomposition into
CY-category summands.
\item Bare $\kappa_\bullet$ without subscript is disallowed; every
$\kappa_\bullet$ above carries its subscript-scope declaration
(one of $\kappa_{\mathrm{cat}}$, $\kappa_{\mathrm{ch}}$,
$\kappa_{\mathrm{BKM}}$, $\kappa_{\mathrm{fiber}}$).
\end{enumerate}
\end{proposition}

\section{Open questions}
\label{sec:open}

\begin{enumerate}
\item \emph{Master $\kappa_{\mathrm{mast}}$?} The naked sum
$\kappa_{\mathrm{cat}} + \kappa_{\mathrm{ch}}^{\mathrm{Heis}}
+ \kappa_{\mathrm{BKM}} + \kappa_{\mathrm{fiber}}
= 0 + 3 + 5 + 24 = 32$ has no intrinsic meaning; the weighted
master
$\kappa_{\mathrm{mast}}(\mathcal A_X) := \sum_i \varepsilon_i
\kappa_{(i)}(\Phi^{(i)}(\mathcal A_X))$
with $\varepsilon_i$ the Kontsevich--Soibelman wall-crossing
weights recovers, at $\mathcal A_{K3 \times E}$, the Borcherds
weight $\kappa_{\mathrm{BKM}}(\Phi_{10}) = c_{10}(0)/2 = 10$.
The pattern of $\varepsilon_i$ pinning is conjectural; the
chain-level verification of the weighted identity across
$N \in \{1, 2, 3, 4, 6\}$ is open.
\item \emph{Pentagon colimit conjecture.} CY-C predicts that the
five non-source routes R$1$, R$3$, R$4$, R$5$, R$6$ assemble into
a pentagon colimit over named intertwiners
$\beta_{13}, \beta_{34}, \beta_{45}, \beta_{56}, \beta_{61}$
through $R_2$ Borcherds as source, with the colimit object
the common chiral-group target $G(K3 \times E)$. Still
conjectural (Vol III CY-C).
\item \emph{Higher-genus extension.} The Schottky--Igusa
obstruction forbids holomorphic continuation of the Gritsenko--Nikulin
tower past $g = 2$; the Zwegers--Dabholkar--Murthy--Zagier
mock-modular completion $\widetilde F_g$ at $g \geq 3$
has mock-modular depth
$d_{\mathrm{mock}}(\widetilde F_g) = \mathrm{codim}_{\mathcal A_g}(\mathcal J_g)$,
equalling $\binom{g-2}{2}$ at $g \geq 4$ (genus $3$ is borderline
with $d_{\mathrm{mock}} = 0$ by Oort--Ueno birational
surjectivity of the Torelli map). Whether this extends to a
higher-genus seventh / eighth route is open.
\item \emph{$(\infty, 1)$-functoriality of $\Phi^{\mathrm{gen}}$.}
Theorem~\ref{thm:stratified-unification} produces
$\Phi^{\mathrm{gen}}$ as an $(\infty, 2)$-stratified functor;
promoting its individual $\Phi^{(i)}$ restrictions to
$(\infty, 1)$-functors requires wall-crossing coherence along
each stratum, Kontsevich--Soibelman style. Open at the
full-theorem level.
\item \emph{The $32$-weighted identity.} Is there a physical
interpretation of the naked sum $32$ in type IIA/M-duality?
$32 = 24 + 8$ (K3 fibre rank plus $c = 8$ denominator) and
$32 = 5 + 27$ (BKM weight plus $27$ fermionic ADE)
admit coincidental arithmetic decompositions. Their
physical content is open.
\end{enumerate}

\appendix

\section{Rectification log}

\noindent Five attack-heal cycles, in order, resolving
potential collapses among the six routes.

\paragraph{Cycle 1: Mukai vs fibre-rank (R$1$ vs R$6$).}
\emph{Attack.} Both land on lattices of rank $24$; same
construction under different names? \emph{Heal.} Different
sublattices, $\mathrm{II}_{3,19}$ vs $\mathrm{II}_{4,20} \oplus
\mathrm{II}_{1,1}(-1)$, giving distinct $\kappa_\bullet$
($\kappa_{\mathrm{cat}}(K3) = 2$ vs
$\kappa_{\mathrm{fiber}} = 24$). \emph{Survived}:
\S\ref{subsec:attack-1}.

\paragraph{Cycle 2: Igusa squaring vs BKM (R$3$ vs R$4$).}
\emph{Attack.} $\Phi_{10} = \Delta_5^2$; one reduces to the
other? \emph{Heal.} Different Narain orbits (light-like $D = 0$
vs imaginary $D > 0$), different Borcherds weights
(light-like shift $\tau(a) = 9$, BKM weight $c_1(0)/2 = 5$).
\emph{Survived}: \S\ref{subsec:attack-2}.

\paragraph{Cycle 3: Monster vs BKM (R$5$ vs R$4$).}
\emph{Attack.} Both BKM superalgebras; both reach Borcherds
lifts? \emph{Heal.} Different root lattices
($\mathrm{II}_{25,1}$ vs $\Lambda^{2,1}_{II}$), different
automorphic inputs ($j$-function coefficients vs $\phi_{0,1}$
Fourier data). \emph{Survived}: \S\ref{subsec:attack-3}.

\paragraph{Cycle 4: CoHA / MO / BFN / Kummer inflate R$2$?}
\emph{Attack.} Four constructions all compute R$2$-like
algebras; true count is $9$? \emph{Heal.} Equivalences within
orbit $(\mathrm{II})$; one route
counted once. \emph{Survived}: \S\ref{subsec:attack-4}.

\paragraph{Cycle 5: Can a single functor unify the four
invariants?} \emph{Attack.} If six constructions produce four
scalars, a surplus suggests unification. \emph{Heal.} No at
the $1$-categorical level (Proposition~\ref{prop:no-single-functor});
yes at the $(\infty, 2)$-stratified level
(Theorem~\ref{thm:stratified-unification}), in which the unifying
functor outputs a sextuple of orbit restrictions whose scalar
shadows collapse onto the tetrad $\{0, 3, 5, 24\}$.
\emph{Survived}: \S\ref{subsec:attack-5}.

\paragraph{Converged content.} The six routes R$1$--R$6$ are
six $O(\Lambda_{\mathrm{Nar}})$-orbits of primitive Mukai
vectors in the Narain lattice $\mathrm{II}_{4,20} \oplus
\mathrm{II}_{1,1}(-1)$, equivalently six U-duality orbits of
cosmic-string types in type IIA on $K3 \times E$, equivalently
six boundary-reduction procedures of 6d hCS on $K3 \times E$,
equivalently six edges out of the basepoint $v_0 = (\mathrm{II}_{3,19},
\Delta_5, \Psi_{\mathrm{Grit}}, \iota_{\Delta_5})$ in a coherent
$6$-simplex $\sigma \in N_6(\mathcal G_{\mathrm{mn}})$. Four of
them produce $\{0, 3, 5, 24\}$; the remaining two (R$3$, R$5$)
produce the light-like $\eta^9$-shift $9$ and the Monster /
Niemeier invariants ($\kappa_{\mathrm{ch}}(V^\natural) = 12$;
$N$-twisted Fourier data). The unifying functor
$\Phi^{\mathrm{gen}}$ exists at $(\infty, 2)$ but not at the
$1$-categorical level. Only R$2$ uses the CY-to-chiral functor
$\Phi$; the other five produce their chiral algebras by
independent means.

\end{document}
